
résolu de
façon satisfaisante la question soulevée au n\textsuperscript{o} 1. Comme support pour la formule (4.5) dans @@M0@@ lui-même, signalons le fait qu'il est clair à priori que la restriction du premier membre à @@M1@@ est nulle, donc, @@M2@@ étant supposé irréductible, le premier membre doit être à priori un multiple entier de @@M3@@:
@@M4@@
Cela prouve déjà la formule exacte (4.5) lorsque @@M5@@ n'est pas un élément de torsion de @@M6@@. Il est immédiat que tel est le cas si @@M7@@ est projectif (et non seulement quasi-projectif) sur le corps @@M8@@, en notant que si @@M9@@ est une polarisation de @@M10@@, alors @@M11@@; cela prouve la formule voulue dans ce cas. Il en résulte que cette formule est encore valable chaque fois que @@M12@@ est isomorphe à un sous-schéma ouvert d'un schéma projectif lisse sur @@M13@@, ce qui sera le cas lorsqu'on dispose d'une solution convenable du problème de résolution des singularités. Il en est ainsi en particulier si @@M14@@ est de caractéristique nulle, grâce à Hironaka [11].

La validité de la formule exacte (4.5) donc de (4.4) impliquerait que le noyau de @@M15@@ est annulé par @@M16@@, et en particulier \uline{que} @@M17@@ \uline{est un isomorphisme}. (Que @@M18@@ soit un isomorphisme résulte déjà immédiatement de X 5.3.2). Sauf erreur, le fait que @@M19@@ soit un isomorphisme n'est pas prouvé à l'heure actuelle, même dans le cas particulier où @@M20@@. C'est hâtivement en tous cas que le rédacteur avait affirmé la validité de (4.4) et (4.5) dans [9, n\textsuperscript{o}~4, 3\textsuperscript{o}].

\textbf{4.4.} A l'heure actuelle, le principal avantage de @@M21@@ (ou de @@M22@@) sur l'anneau de Chow @@M23@@ est le fait qu'il se prête mieux aux calculs explicites de certaines formules d'intersection entre cycles, dans le cas d'intersections ``excédentaires'' i.e. de dimension trop grande. La définition de la classe d'intersections, à l'aide du lemme de Chow permettant de bouger suffisamment un cycle dans sa classe, se prête en effet malaisément au calcul. Un exemple frappant à cet égard est la formule fondamentale suivante [Exp. VII 2.7], pour une immersion fermée @@M24@@ de schémas quasi-projectifs lisses sur @@M25@@:
@@M26@@
formule valable dans @@M27@@, @@M28@@ étant dans @@M29@@, @@M30@@ le faisceau conormal 
